\section*{Capítulo \ref{ch:b3:heterocycles} --- Química dos heterociclos}

\begin{solution}{exo:b3:heterocycles:1}
Oxazol: O dá dois (um par isolado no sistema $\pi$, um no plano), N um, os
três carbonos três: 6. Tiazol: o mesmo com S: 6. Pirimidina: um de cada
um dos seis átomos, os dois pares isolados dos nitrogênios no plano: 6. Purina: nove átomos,
o nitrogênio N--H dando dois e os outros um cada: 10, um número de Hückel
($4n + 2$, $n = 2$).
\end{solution}

\begin{solution}{exo:b3:heterocycles:2}
Piperidina (11.1) $>$ DMAP (9.6) $>$ imidazol (7.0) $>$ piridina (5.2) $\gg$ pirrol (cerca de $-4$,
protonado no carbono).
\end{solution}

\begin{solution}{exo:b3:heterocycles:3}
2-Bromopirrol (o pirrol é tão reativo que se polibroma facilmente);
2-bromofurano; 3-bromoindol.
\end{solution}

\begin{solution}{exo:b3:heterocycles:4}
2,5-Dimetilfurano; 1,2,5-trimetilpirrol; 2,5-dimetiltiofeno.
\end{solution}

\begin{solution}{exo:b3:heterocycles:5}
2-Metoxipiridina. A adição do metóxido em C2 põe a carga negativa do
\omterm{def:b3:heterocycles:snar}{complexo de Meisenheimer} no nitrogênio; em C3, ela só pode se distribuir sobre carbonos, de modo que o
intermediário, e o estado de transição que o precede, têm energia muito mais alta.
\end{solution}

\begin{solution}{exo:b3:heterocycles:6}
2-Aminopiridina (como seu sal de sódio, depois protonada no isolamento). O íon amideto
se adiciona em C2, dando um aduto $\sigma$ aniônico com H e $\ce{NH2}$ em C2 e a carga
negativa no nitrogênio do anel; ele perde um hidreto, que toma um próton do
grupo amino e sai como $\ce{H2}$.
\end{solution}

\begin{solution}{exo:b3:heterocycles:7}
O oxigênio do N-óxido devolve um par isolado ao anel: as estruturas de ressonância
têm carga negativa em C2, C4 e C6. Os eletrófilos atacam C4 (C2
fica perto do nitrogênio carregado positivamente). O oxigênio é depois removido com
tricloreto de fósforo, dando a 4-nitropiridina.
\end{solution}

\begin{solution}{exo:b3:heterocycles:8}
2-Hidroxipiridina (um OH em C2) e 2-piridona (N--H e C=O). A piridona conserva seis
elétrons $\pi$ no anel: sua estrutura de ressonância zwitteriônica, com $\mathrm{O^-}$ e um
$\mathrm{N^+}$ que compartilha seu par isolado, é um piridínio-olato; ela é uma amida cujo
par isolado do nitrogênio completa o sexteto, e as fortes ligações de hidrogênio de C=O e de N--H
a favorecem em água.
\end{solution}

\begin{solution}{exo:b3:heterocycles:9}
2-Nitrobenzaldeído, dois equivalentes de acetoacetato de metila e amônia: o
bloqueador dos canais de cálcio nifedipino.
\end{solution}

\begin{solution}{exo:b3:heterocycles:10}
A eno-hidrazina voltada para o grupo $\ce{CH2}$ (interna, mais substituída) dá o
2,3-dimetilindol; a voltada para o grupo metila (terminal) dá o 2-etilindol.
Os ácidos fracos dão sobretudo o primeiro, pela eno-hidrazina mais estável;
ácidos mais fortes aumentam a proporção de 2-etilindol.
\end{solution}

\begin{solution}{exo:b3:heterocycles:11}
Na $\mathrm S_\mathrm NAr$, a etapa determinante da velocidade é a adição, que não rompe a ligação C--X;
o flúor, mais eletronegativo, estabiliza mais o estado de transição aniônico
e o \omterm{def:b3:heterocycles:snar}{complexo de Meisenheimer}. Na $\mathrm S_\mathrm N2$, a ligação C--X se rompe no
estado de transição, e a fraca ligação C--I é a que sai mais depressa.
\end{solution}

\begin{solution}{exo:b3:heterocycles:12}
Piridina: três sinais na razão 2\,:\,1\,:\,2, o primeiro perto de \qty{8.6}{ppm}. Pirrol: dois sinais
iguais em 6.7 e \qty{6.2}{ppm} e um N--H largo perto de \qty{8.1}{ppm}. Furano: dois sinais iguais em
7.4 e \qty{6.4}{ppm}, sem N--H. Tiofeno: dois sinais iguais em 7.33 e \qty{7.12}{ppm}, próximos
um do outro, entre os do furano.
\end{solution}

\begin{solution}{pb:b3:heterocycles:1}
\textbf{1.} $\ce{C6H5NHNH2 + (CH3)2CO -> C6H5NHN=C(CH3)2 + H2O}$.

\textbf{2.} O nitrogênio amínico se adiciona ao carbono carbonílico, depois se perde água, como na
formação de uma imina; o ácido ativa o grupo carbonila e transforma o OH do aduto
em grupo de saída.

\textbf{3.} O $\ce{NH2}$ terminal: o par isolado do outro nitrogênio está conjugado com o
anel fenila e é mais impedido.

\textbf{4.} $\qty{108.0}{g/mol}$; $\qty{10.0}{g}/\qty{108.0}{g/mol} = \qty{0.0926}{mol}$.

\textbf{5.} A hidrazona $\ce{C9H12N2}$, \qty{148.0}{g/mol}: $\qty{0.0926}{mol} \times \qty{148.0}{g/mol} = \qty{13.7}{g}$.

\textbf{6.} $\ce{C6H5-NH-NH-C(CH3)=CH2}$.

\textbf{7.} O carbono \textit{orto} do anel, o carbono \textit{ipso}, os dois nitrogênios, o carbono da hidrazona
e o $\ce{CH2}$ terminal.

\textbf{8.} A ligação N--N se rompe; forma-se uma ligação C--C entre o carbono \textit{orto} e o $\ce{CH2}$.

\textbf{9.} Seis elétrons, todos os componentes \omterm{def:b3:pericyclic:faciality}{suprafaciais}: termicamente permitida.

\textbf{10.} Como ácido de Lewis, ele catalisa a tautomerização, liga um nitrogênio e
enfraquece a ligação N--N, e ajuda a saída da amônia.

\textbf{11.} O carbono \textit{orto} tornou-se $sp^3$, com um hidrogênio: a perda desse próton
restaura o anel benzênico, uma forte força motriz.

\textbf{12.} O $\ce{NH2}$ anilínico formado ataca o carbono da imina, fechando o
anel de cinco membros (uma 2-aminoindolina).

\textbf{13.} A amônia; a formação do anel pirrólico aromático do indol a impulsiona.

\textbf{14.} $\ce{C9H12N2 -> C9H9N + NH3}$.

\textbf{15.} \qty{131.0}{g/mol}.

\textbf{16.} $\qty{0.0926}{mol} \times \qty{131.0}{g/mol} = \qty{12.1}{g}$.

\textbf{17.} Em C3, a posição mais nucleofílica do indol, aqui livre; o ataque nessa posição
mantém intacto o anel benzênico no intermediário.

\textbf{18.} $\ce{C6H5NHNH-C(=CH2)-CH2CH3}$ (terminal) e $\ce{C6H5NHNH-C(CH3)=CH-CH3}$ (interna).

\textbf{19.} 2-Etilindol e 2,3-dimetilindol.

\textbf{20.} A interna é mais substituída; os ácidos fortes aumentam a parcela da
terminal.

\textbf{21.} O 2,3-dimetilindol, pela eno-hidrazina mais estável e mais
substituída.

\textbf{22.} O acetaldeído se condensa consigo mesmo e a rota da eno-hidrazina dá baixos
rendimentos; o indol é preparado, em vez disso, a partir da hidrazona do ácido pirúvico, que dá o
ácido indol-2-carboxílico, depois descarboxilado.

\textbf{23.} É preciso separar uma mistura de regioisômeros, o que custa rendimento; prefere-se uma
cetona simétrica ou uma rota que fixe a regioquímica.

\textbf{24.} 10.0 g de fenil-hidrazina podem dar no máximo \qty{12.1}{g} de 2-metilindol.
\end{solution}
